\section{Introduction}
\label{sec:intro}

Deep segmentation models match expert radiologists on brain-tumor benchmarks, yet clinical
deployment lags because their decisions cannot be verifiably explained. The dominant remedy is
post-hoc attribution~\cite{selvaraju2017gradcam,sundararajan2017ig,lundberg2017shap}, which
produces a saliency map \emph{after} the prediction and therefore approximates the mechanism
rather than reporting it. Rudin~\cite{rudin2019stop} argues for the alternative: build models
whose decisions are interpretable by construction. Concept Bottleneck Models
(CBMs)~\cite{koh2020cbm} realize this for classification by forcing the prediction through a
human-readable concept layer, so that the concept activation \emph{is} the explanation.

Extending that idea to dense prediction is attractive, and---as we show---treacherous. We build
the natural extension: replace the encoder bottleneck of a segmentation network with a router
that assigns every token to one of $K$ clinically labeled experts $\{$Background, Necrotic core,
Edema, Enhancing tumor$\}$. The routing distribution is a per-region, concept-labeled map
computed in a single forward pass with no post-hoc step. It looks exactly like an intrinsic
explanation should.

\paragraph{It aligns beautifully, and it does nothing.}
Trained end-to-end, the routing tracks the clinical concepts closely ($\casfg$ up to $0.96$) and
discriminates them far better than any post-hoc method ($0.88$ vs.\ $0.74$ for the strongest
baseline). These are the metrics this literature reports, and by both the module is a success.
But when we \emph{intervene}---force the tokens the router assigned to one concept through a
different concept's expert, changing nothing else---the prediction does not move. Only $1.1\%$ of
affected voxels change label, and zeroing the entire bottleneck costs $1.5$ Dice on edema.

The failure is invisible to alignment metrics \emph{by construction}. $\casfg$ measures
correlation between the routing and the ground truth, and correlation is perfectly compatible
with a bottleneck the model never consults. Worse, the alignment loss trains precisely what
$\casfg$ scores, so a high value confirms that the router learned its supervision---not that the
network uses it.

\paragraph{Why it happens.}
We rule out the obvious explanation: the experts have \emph{not} collapsed to a shared function.
They displace tokens by $59$--$74\%$ of their norm and differ pairwise by $0.93$. The cause is
architectural, and it is two independent leaks that encoder--decoder designs invite. The expert
computes $h + c_k(h)$, so the residual carries the full feature vector to the decoder whichever
expert fires---routing modulates the representation but never gates it. And the decoder's skip
connections carry encoder features directly to the output, so it can decide what tissue it is
looking at without consulting the routing. Either leak alone suffices, and neither is visible in
$\casfg$.

\paragraph{A control that reframes the baseline.}
We also find that a linear probe on an encoder trained with \emph{no} concept supervision already
separates the concepts at $0.81$---above every post-hoc method we test. Concept-discriminability
gains over saliency baselines are therefore much weaker evidence for a concept module than the
literature treats them as.

\paragraph{Closing the leaks.}
Attempts to force dependence by \emph{removing} the decoder's alternatives fail: a scalar gate on
the skip connections is compensated away during training, and removing the expert residual
backfires, collapsing pairwise expert divergence from $0.93$ to $0.16$ because every expert must
then independently reconstruct the full signal. Making the routing an \emph{additive term of the
output} succeeds, because causal influence becomes arithmetic rather than something the optimizer
can route around. Causal control rises from $0.011$ to $0.636$, removing the routing term costs
up to $99\%$ Dice, and the price is ${\approx}1$ Dice point.

\noindent\textbf{Contributions.}
\begin{enumerate}\itemsep2pt
\item \textbf{A diagnostic protocol for concept bottlenecks} (Sec.~\ref{sec:protocol}): causal
intervention, bottleneck ablation, and a linear-probe control. Three cheap tests that separate a
bottleneck the model uses from one it ignores. Alignment metrics cannot.
\item \textbf{Evidence that the natural design fails it} (Sec.~\ref{sec:inert}). A bottleneck with
$\casfg$ up to $0.96$ and best-in-class concept discriminability is causally inert, and we localize
the cause to two specific architectural leaks.
\item \textbf{A repair and its price} (Sec.~\ref{sec:repair}). Additive routing makes the
bottleneck causally decisive---$0.636$ $[0.621,0.651]$ label-flip under intervention against
$0.011$ $[0.009,0.014]$, $p<10^{-28}$---for ${\approx}1$ Dice point, with alignment and calibration
slightly improved.
\end{enumerate}
